% META: {"id": "1SPE-TRIGO-ME-004", "chapitre": "1SPE-TRIGONOMETRIE", "type_objet": "methode", "capacites": ["1SPE-TRIGONOMETRIE-C4"], "capacites_codes": ["C4"], "methodes": ["M4"], "mode_creation": "ex_nihilo", "sources_inspiration": [], "status": "needs_review"}
% BEGIN-VERIFY
% from sympy import *
% x = symbols('x')
% # Exemple redige : cos(x) = sqrt(2)/2 sur ]-pi ; pi]
% sols = solveset(Eq(cos(x), sqrt(2)/2), x, Interval.Lopen(-pi, pi))
% assert sols == FiniteSet(-pi/4, pi/4)
% # Verification par valeurs
% assert simplify(cos(pi/4) - sqrt(2)/2) == 0
% assert simplify(cos(-pi/4) - sqrt(2)/2) == 0
% # Piege : sin(x) = 1/2 sur ]-pi ; pi] a pour solutions pi/6 et 5*pi/6 (pas -pi/6)
% sols2 = solveset(Eq(sin(x), Rational(1,2)), x, Interval.Lopen(-pi, pi))
% assert sols2 == FiniteSet(pi/6, 5*pi/6)
% END-VERIFY
\begin{fichemethode}{M4}{Resoudre cos(x) = a ou sin(x) = a sur un intervalle}

\quandUtiliser{%
  Utiliser cette methode lorsque l'enonce demande de resoudre une equation
  $\cos(x) = a$ ou $\sin(x) = a$ (ou une inequation associee) sur un
  intervalle donne, le plus souvent $]-\pi\,;\,\pi]$ ou $[0\,;\,2\pi[$.
}

\pasApas{%
  \begin{enumerate}
    \item Verifier que $a \in [-1\,;\,1]$ : sinon l'equation n'a aucune solution.
    \item Reconnaitre une valeur remarquable : ecrire $a$ sous la forme
      $\cos(\alpha)$ (ou $\sin(\alpha)$) avec $\alpha$ un angle connu
      ($0$, $\frac{\pi}{6}$, $\frac{\pi}{4}$, $\frac{\pi}{3}$, $\frac{\pi}{2}$...).
    \item Ecrire TOUTES les solutions :
      \begin{itemize}
        \item $\cos(x) = \cos(\alpha) \iff x = \alpha + 2k\pi$ ou $x = -\alpha + 2k\pi$, $k \in \mathbb{Z}$ ;
        \item $\sin(x) = \sin(\alpha) \iff x = \alpha + 2k\pi$ ou $x = \pi - \alpha + 2k\pi$, $k \in \mathbb{Z}$.
      \end{itemize}
    \item Selectionner les solutions appartenant a l'intervalle demande en
      choisissant les valeurs de $k$ convenables.
    \item Controler chaque solution sur le cercle trigonometrique.
  \end{enumerate}
}

\exempleRedige{%
  Resoudre dans $]-\pi\,;\,\pi]$ l'equation $\cos(x) = \dfrac{\sqrt{2}}{2}$.

  \medskip
  \commentaireMarge{$\frac{\sqrt{2}}{2} = \cos\frac{\pi}{4}$ : valeur remarquable.}%
  On reconnait $\dfrac{\sqrt{2}}{2} = \cos\left(\dfrac{\pi}{4}\right)$, donc
  l'equation s'ecrit $\cos(x) = \cos\left(\dfrac{\pi}{4}\right)$.

  \medskip
  Les solutions reelles sont $x = \dfrac{\pi}{4} + 2k\pi$ ou
  $x = -\dfrac{\pi}{4} + 2k\pi$, avec $k \in \mathbb{Z}$.

  \medskip
  \commentaireMarge{On garde les solutions dans $]-\pi;\pi]$.}%
  Dans $]-\pi\,;\,\pi]$, seules conviennent ($k = 0$) :
  \[
    S = \left\{ -\dfrac{\pi}{4}\,;\ \dfrac{\pi}{4} \right\}.
  \]
}

\pieges{%
  \begin{itemize}
    \item \textbf{Piege 1.} Oublier la seconde famille de solutions :
      $\cos(x) = \cos(\alpha)$ donne $x = \alpha$ \emph{et} $x = -\alpha$
      (a $2k\pi$ pres) ; pour le sinus c'est $\alpha$ \emph{et} $\pi - \alpha$.
    \item \textbf{Piege 2.} Confondre les deux regles : pour
      $\sin(x) = \frac{1}{2}$ sur $]-\pi\,;\,\pi]$, les solutions sont
      $\frac{\pi}{6}$ et $\frac{5\pi}{6}$ — pas $-\frac{\pi}{6}$.
    \item \textbf{Piege 3.} Ne pas ajuster $k$ a l'intervalle demande :
      sur $[0\,;\,2\pi[$, la solution negative $-\frac{\pi}{4}$ devient
      $-\frac{\pi}{4} + 2\pi = \frac{7\pi}{4}$.
  \end{itemize}
}

\verifier{%
  Remplacer chaque solution dans l'equation de depart, et verifier sur le
  cercle trigonometrique que le nombre de points d'intersection avec la
  droite $x = a$ (ou $y = a$) correspond au nombre de solutions trouvees
  dans l'intervalle.
}

\sEntrainer{\refExos{M4}}

\end{fichemethode}
